\chapter{RESUMEN DEL PROYECTO}

\section{Contexto y justificación}

La contaminación por microplásticos es un desafío medioambiental global, y la espectroscopía Raman es hoy la técnica habitual para identificar el polímero de una partícula por comparación con bibliotecas espectrales (\textit{matching}). La calidad del espectro medido depende directamente del tiempo de integración: un tiempo corto, deseable para aumentar el rendimiento de muestreo, produce espectros ruidosos y más difíciles de identificar. El comportamiento de las distintas métricas de similitud frente a esa degradación rara vez se compara de forma sistemática.

\section{Planteamiento del problema}

¿Cómo se degrada el \textit{matching} de un espectro Raman de microplástico contra una base de datos de referencia a medida que empeora la calidad de la señal, y qué métrica de similitud es más robusta frente a esa degradación? El proyecto tiene dos mitades: caracterizar empíricamente cómo se comportan siete métricas frente al ruido, y construir el \textit{pipeline} que lo hace posible, con preprocesado, \textit{matching} y análisis estadístico, reproducible y paralelizable con Apache Spark.

\section{Objetivos del proyecto}

El objetivo general es evaluar a gran escala la robustez de distintas métricas de similitud espectral frente a la degradación de la calidad del espectro Raman, implementando un \textit{pipeline} reproducible y paralelizable. Los objetivos específicos se detallan en el Capítulo~3: validar el preprocesado contra el protocolo de referencia, implementar las siete métricas, diseñar y paralelizar la matriz experimental, medir la robustez con el punto de ruptura y un análisis estadístico, y generar el entregable visual.

\section{Resultados obtenidos}

El \textit{pipeline} reproduce el protocolo de referencia con diferencia numérica nula y el experimento se ha ejecutado completo sobre cinco polímeros y la biblioteca pública \emph{Hagelskjær Microplastic Solution Library~1.1}, superando la validación numérica acordada (1885 comparaciones alineadas). Pearson aguanta mejor la acumulación real (61,5~\% de aciertos); con la degradación paramétrica el mejor grupo es otro, el de coseno, euclídea, HQI y SAM (87,5~\%). Manhattan y Spearman quedan por debajo en los dos ejes, y ninguna diferencia por pares sobrevive la corrección de Holm. Se detectó además que la biblioteca de referencia no cubre dos de los cinco polímeros medidos.

\section{Estructura de la memoria}

La memoria se organiza en nueve capítulos. El Capítulo~2 presenta los antecedentes y el estado del arte; el Capítulo~3 detalla los objetivos; el Capítulo~4, núcleo del documento, describe la planificación, el \textit{pipeline}, las decisiones técnicas y los resultados; el Capítulo~5 discute limitaciones y adaptaciones; los Capítulos~6 y~7 recogen conclusiones y futuras líneas de trabajo; los Capítulos~8 y~9, las referencias bibliográficas y los anexos técnicos.
